\section*{Chapter \ref{ch:dv:dividends-borrow-and-forwards} --- Dividends, Borrow and Forwards}

\begin{solution}{exo:dv:dividends-borrow-and-forwards:1}
$(100-e^{-0.0075}-e^{-0.0225})e^{0.03}=101.0152$.
\end{solution}

\begin{solution}{exo:dv:dividends-borrow-and-forwards:2}
Slope $(-2.83-7.02)/10=-0.985$: the discount factor is 0.985. Then $7.02=0.985(F-95)$ gives
$F=102.1269$.
\end{solution}

\begin{solution}{exo:dv:dividends-borrow-and-forwards:3}
$\bar D=4e^{-0.015}=3.9404$, $S/(S-\bar D)=1.0410$;
$\sigma_{\mathrm{eff}}=0.25\sqrt{1.0410^2\times0.5+0.5}=0.2552$.
\end{solution}

\begin{solution}{exo:dv:dividends-borrow-and-forwards:4}
Each side can be exercised early, so $C-P$ is only bounded:
$S-D-K\le C-P\le S-Ke^{-rT}$ for American options on a dividend payer. Before an ex-date
the call carries an early-exercise premium (it may be exercised to capture the
dividend), so $C-P$ is too high and the forward read from the European relation is too
high: the \omterm{def:dv:dividends-borrow-and-forwards:implied}{implied dividend} comes out too low.
\end{solution}

\begin{solution}{exo:dv:dividends-borrow-and-forwards:5}
Cash: $F=(S-4e^{-0.015})e^{0.03}$, so $\partial F/\partial S=e^{0.03}=1.0305$.
Proportional: $F=0.96\,Se^{0.03}$, so $\partial F/\partial S=0.9892$. The same dividend
gives two hedges 4\% apart.
\end{solution}

\begin{solution}{exo:dv:dividends-borrow-and-forwards:6}
The forward rises from 100.3611 to 101.1038 and the call from 12.7558 to 13.1475: $+0.3917$,
about 3\% of its value, for a 10\% cut in dividends.
\end{solution}

\begin{solution}{exo:dv:dividends-borrow-and-forwards:7}
Forward 100.1132, discount factor 0.97033; with the known dividends (present value 2.3661)
the implied borrow is 0.504\%, against 0.5\%: the regression error, about a cent on the
forward, is 0.004\% of borrow.
\end{solution}

\begin{solution}{exo:dv:dividends-borrow-and-forwards:8}
Both volatilities are computed on the wrong forward. On the forward the options price,
48.74, the call and the put at 50 both imply 60\%: there is no volatility difference to
capture. The trade is a reversal in disguise: selling puts and buying calls, hedged by
shorting shares, needs a borrow at 35\%, which costs what the ``edge'' appears to pay.
\end{solution}

\begin{solution}{pb:dv:dividends-borrow-and-forwards:1}
\textbf{1.} 3.73, 1.23, $-1.26$, $-3.75$, $-6.24$.
\textbf{2.} Discount factor 0.99680; forward 48.74.
\textbf{3.} 3.90\%; the quotes are rounded to the cent, and a cent of error on a slope
measured over ten points of strike moves the discount factor by $10^{-3}$.
\textbf{4.} 1.42.
\textbf{5.} 35.0\% a year.
\textbf{6.} 50.16.
\textbf{7.} Call 3.50, put 3.34.
\textbf{8.} The call by $+0.72$ (26\% of its price), the put by $-0.70$.
\textbf{9.} Call 47.3\%, put 72.3\%.
\textbf{10.} Parity ties $C-P$ to the forward; with the naive forward the call's and the
put's volatilities must differ to fit it. On the forward 48.74 both are 60\%.
\textbf{11.} $51.26-50e^{-0.04\times30/365}$: 1.42 today.
\textbf{12.} $50\times0.35\times30/365=1.44$.
\textbf{13.} No: the apparent profit is the borrow cost, to within the rounding of the quotes.
\textbf{14.} Sell the share and buy it back synthetically (buy the call, sell the put):
this brings in 1.42 for the month, the same as lending the share at 35\%. Holders who
cannot lend (a fund whose custodian does not lend) earn the fee this way.
\textbf{15.} Buy puts or put spreads, or sell calls: the options carry the borrow cost in
their prices, so the view is expensive but expressible.
\textbf{16.} The forward jumps from 48.74 to 49.96: the 50 call rises and the put falls by
about 0.6 each, a move unrelated to volatility or spot.
\textbf{17.} The borrow fee acts like a dividend yield: an American call may be worth
exercising early to own the share and lend it at the fee, which adds an early-exercise
premium to the calls and blurs parity further.
\textbf{18.} The forward implied by the options themselves, updated with the borrow market,
not spot times the financing rate.
\textbf{19.} 35.0\% a year; the naive price of the 50 call is 0.72 too high.
\textbf{20.} The forward, and through it the cost of being short the share.
\end{solution}

\begin{solution}{iq:dv:dividends-borrow-and-forwards:1}
Put--call parity: $C-P=P(0,T)(F-K)$. Regress $C-P$ on $K$ over European quotes near the
money; the slope is $-P(0,T)$, the intercept $P(0,T)F$. With American options, use
out-of-the-money quotes where early exercise is negligible or remove its premium first.
\iqlookfor{parity as a linear regression across strikes.}
\end{solution}

\begin{solution}{iq:dv:dividends-borrow-and-forwards:2}
The forward is below the spot: a large dividend before expiry, or a high borrow fee (a
hard-to-borrow share). Both appear in parity as carry; at the forward the two options are
worth the same.
\iqlookfor{dividends and borrow through the forward.}
\end{solution}

\begin{solution}{iq:dv:dividends-borrow-and-forwards:3}
Cash for the near dividends, already declared; proportional (or a blend) for far ones,
because a company's dividend falls when its share falls. The hedge changes because the
forward's sensitivity to spot is $1/P(0,T)$ with cash dividends and $F/S$ with
proportional ones; for a five-year option the gap is of the order of the dividend yield
times five years.
\iqlookfor{the dynamics, not only the forward level.}
\end{solution}

\begin{solution}{iq:dv:dividends-borrow-and-forwards:4}
The sensitivity of positions to expected dividends. Issuers of equity-linked
structured products, whose hedges collect dividends they have promised to note holders,
are long; they sell \omterm{def:dv:dividends-borrow-and-forwards:risk}{dividend swaps} and futures. In 2020 supervisors asked banks to suspend
dividends and index dividends traded far below models, and issuers' hedging sales added
to the fall.
\iqlookfor{who carries it, and why the 2020 fall was amplified.}
\end{solution}

\begin{solution}{iq:dv:dividends-borrow-and-forwards:5}
Options: escrowed (a tree on $S-\mathrm{PV}(D)$, adding the PV back for exercise), or a
drop of the node values by $D$ at the ex-date. The escrowed tree understates the
volatility before the ex-date unless its volatility is adjusted; dropping by $D$ breaks
recombination (the tree must be rebuilt after the ex-date or interpolated), and a
proportional drop keeps recombination but is wrong for a cash dividend. Test against a
European reference with the same forward.
\iqlookfor{the recombination problem and the volatility mismatch.}
\end{solution}

\begin{solution}{iq:dv:dividends-borrow-and-forwards:6}
The spot model: its volatility applies to the whole share including the part paid out, so
before the ex-date the forward moves more. For a one-year option with a 4\% dividend at
six months, about half a volatility point, 0.2 of premium at the money; the adjustment is
$\sigma^2(S/(S-\bar D))^2t_1+\sigma^2(T-t_1)$.
\iqlookfor{what the volatility is the volatility of.}
\end{solution}
